\documentclass[11pt]{article}
\usepackage[no-math]{fontspec}
\usepackage[french,english]{babel}
\usepackage{amsmath}
\usepackage{amssymb}
\usepackage{amsthm}
\usepackage{mathrsfs}
\usepackage{tikz-cd}
\usepackage{enumitem}

\frenchsetup{AutoSpacePunctuation=false}

\newcommand{\printedpage}[1]{\marginpar{\footnotesize p.~#1}}
\newcommand{\msp}[1]{\marginpar{\footnotesize [#1]}}
\newcommand{\deq}{\overset{\text{def}}{=}}

\begin{document}

% ===================== page 1 of the PDF (title page) =====================
\begin{center}
\printedpage{[no number]}% no page number printed on this page
% head: (none)

\vspace*{6cm}

{\Large\textsf{THE DERIVATORS}}

\bigskip

\textbf{ALEXANDRE GROTHENDIECK}

\vspace{2cm}

\textbf{Chapter VI}

\medskip

\textbf{HOT}
\end{center}

\clearpage

% ===================== page 2 of the PDF (editorial note) =====================
\vspace*{6cm}

\noindent\printedpage{[no number]}% no page number printed on this page
% head: (none)
This text was deciphered and transcribed in \LaTeXe{} by M. Künzer. It was edited by M. Künzer, J. Malgoire, and G. Maltsiniotis. The transcription is as faithful as possible to the manuscript. For the few obvious corrections, or rare comments by the editors, as well as for the original numbering of the manuscript pages, typewriter characters \texttt{[typewriter]} in brackets are used. A question mark in brackets means that one is not sure of the preceding word.

\noindent This edition is provisional. Remarks, comments and corrections are welcome. Send a message to:

\noindent\textbf{[e-mail address removed]}

\hfill G. Maltsiniotis\hspace*{2cm}

\clearpage

% ===================== printed page 1 =====================
\noindent\printedpage{1}
% head: A. Grothendieck : Dérivateurs, Ch. VI, Édition des Universités de Montpellier II et Paris VII
% foot: Version 9 novembre 2001
\msp{page 1}

\section*{\normalsize The derivator \textmd{HOT}}

\begin{flushright}
14.11.90
\end{flushright}

\begin{flushright}
(The first sheets, from two weeks ago, have been ditched, overtaken by VII)
\end{flushright}

\noindent I define Hot starting from the category of models
\[
\mathcal{M} = \mathrm{Cat},\ W \subseteq \mathrm{Fl}(\mathrm{Cat}),
\]
$W$ being in principle the Hot-equivalences. But as far as possible I am going to axiomatize more, assuming given a more general set $W$, called ``basic localizer'' (basic localizer), e.g. $W = W_{\mathbf{D}}$, where $\mathbf{D}$ is an already given derivator (but assuming that its domain is all of (Cat)). We are going to note as we go along which properties of $W$ we will need. For purposes of later reference, I am going to recall here those I isolated in Pursuing Stacks.

\noindent\textbf{L1}\ \ Saturation of $W$

\begin{itemize}[leftmargin=2.5em]
\item[a)] \textit{The identities} \texttt{[are]} \textit{in $W$.}
\item[b)] \textit{If two among $gf$, $f$, $g$} \texttt{[are]} \textit{in $W$, the third also, $X \xrightarrow{f} Y \xrightarrow{g} Z$.}
\item[c)] \textit{If $gf$ and $fg$ are in $W$, then $f$ and $g$ too, $X \overset{f}{\underset{g}{\rightleftarrows}} Y$ .}
\end{itemize}

\noindent\textbf{L2}\ \ $\mathbf{\Delta}^1$ \textit{is universally $W$-aspheric, i.e. $\mathbf{\Delta}^1 \times X \longrightarrow X$ is in $W$ for every $X$ in} Cat. This implies

\noindent\textbf{L2 bis}\ \ $\mathbf{\Delta}^1$ \textit{is $W$-aspheric.} (This was L3$'$ in Pursuing Stacks.)

\noindent\textbf{L3}\ \ \textit{If $X$ has a final object, $X$ is $W$-aspheric.} (\footnote[1]{There is an L3$'$ dual to L3.})

\noindent I note here \texttt{[some]} strengthened variants.

\noindent\textbf{L3 bis}\ \ \textit{If $X$ is} Hot-\textit{aspheric, it is $W$-aspheric.}

\noindent\textbf{L3 ter}\ \ \textit{If $X$ is} Hot-\textit{aspheric, it is universally $W$-aspheric, i.e. $X \times Y \xrightarrow{\mathrm{pr}_2} Y$ is in $W$ for every $Y$ in} Cat. (\footnote[2]{Suffices to imply L2, L3.})

\clearpage

% ===================== printed page 2 =====================
\noindent\printedpage{2}
% head: A. Grothendieck : Dérivateurs, Ch. VI, Édition des Universités de Montpellier II et Paris VII
% foot: Version 9 novembre 2001
\textbf{L3 quater}\ \ \textit{Every} Hot-\textit{equivalence is in $W$.}

\noindent L3 quater $\Longrightarrow$ L3 ter $\Longrightarrow$ L3 bis $\Longrightarrow$ L3.

\noindent L3 ter $\Longrightarrow$ L2 $\Longrightarrow$ L2 bis.

\noindent\msp{page 2}

\begin{itemize}[leftmargin=2.5em]
\item[\textbf{L4}] \textit{If $f : X \longrightarrow Y$ is $W$-aspheric (i.e. the $X/y \longrightarrow Y/y$ in $W$, i.e. (if we have L3) if the $X/y$ are $W$-aspheric), then $f$ too.} (Localization axiom.) (\footnote[3]{there is an L4$'$ dual to L4.}) \texttt{[``also''} \texttt{= in} $W$\texttt{]}

\item[\textbf{L5}] (Fibration axiom). \textit{If}
\[
\begin{tikzcd}[column sep=small]
X \arrow[rr, "f"] \arrow[dr] & & Y \arrow[dl] \\
& S &
\end{tikzcd}
\]
\textit{with $X$ and $Y$ fibered over $S$ and $f$ cartesian, if for every $s \in S$ $f_s : X_s \longrightarrow Y_s$ is in $W$, $f$ too.}
\end{itemize}

\noindent In Pursuing Stacks I always assumed \fbox{L1, L3, L4} (implying L2) - it was these conditions that defined the notion of ``localizer'' (basic) in Cat. (\footnote[4]{If $W$ \texttt{[is]} a localizer, then contractible categories are $W$-aspheric, homotopisms are in $W$, homotopic maps are equal in (Cat)$W^{-1} \deq \mathrm{Hot}_W$.}) On the other hand, \fbox{L1, L2 bis, L5} imply L1,2,3,4,5.

\noindent\textit{We are going to assume here too that we have at least} L1, L3, L4 (i.e. $W$ a localizer). Applying the construction of I, p.~72-\textbf{75,} we see that $f : X \longrightarrow Y$ is in $W$ if and only if $f^o : X^o \longrightarrow Y^o$ is. So L1, L3, L4 imply L3$'$, L4$'$ (axioms dual to L3, L4).

\noindent There were supplementary axioms

\begin{itemize}[leftmargin=2.5em]
\item[\textbf{La}] $f \in W \Longrightarrow \pi_0(f)$ \textit{bijective.}

\item[\textbf{Lb}] \textit{The canonical localization functor}
\[
\mathrm{Cat} \longrightarrow (\mathrm{Cat})W^{-1} \deq \mathrm{Hot}_W
\]
\textit{commutes with finite sums.}

\item[\textbf{Lc}] \textit{Likewise for finite products.}
\end{itemize}

\clearpage

% ===================== printed page 3 =====================
\noindent\printedpage{3}
% head: A. Grothendieck : Dérivateurs, Ch. VI, Édition des Universités de Montpellier II et Paris VII
% foot: Version 9 novembre 2001
One should all the same isolate concrete conditions on $W$ for the validity of Lb, Lc. The conditions L3 quater and La are written respectively
\[
W_\infty \subseteq W,\ \ W \subseteq W_0.
\]

\noindent\msp{page 3}

\noindent Their conjunction is the very reasonable axiom
\[
W_\infty \subseteq W \subseteq W_0.
\]
Of course, axiom L5 now looks a bit cramped, i.e. feeble. We are going to see very soon what we will need!

\noindent We thus want to define a derivator
\[
\mathrm{HOT}_W
\]
starting from the category of models
\[
(\mathrm{Cat}, W \subseteq \mathrm{Fl}(\mathrm{Cat})).
\]
So to define
\[
\mathrm{HOT}_W(S) \qquad\qquad\qquad (S \in \mathrm{Ob}\,\mathrm{Cat}),
\]
the general recipe is to take the category % [sic] récette -> recette, misspelling of "recette"
\[
\mathcal{M}(S) = \mathrm{Cat}(S) = \underline{\mathrm{Hom}}(S^o, (\mathrm{Cat})),
\]
to introduce in it the set of arrows
\[
\left\{
\begin{aligned}
& W_S \subseteq \mathrm{Fl}(\mathcal{M}(S)) \\
& u \in W_S \Longleftrightarrow u_s \in W \ \ \forall\, s \in \mathrm{Ob}\,S.
\end{aligned}
\right.
\]
We must set
\[
\mathrm{HOT}_W(S) \deq \underline{\mathrm{Hom}}(S^o, \mathrm{Cat})W_S^{-1}.
\]
But the category $\underline{\mathrm{Hom}}(S^o, \mathrm{Cat})$ is interpreted, up to canonical isomorphism, as that
\[
\underline{\mathrm{Fibsc}}\,S
\]
of categories $X$ Cat-fibered over $S$, equipped with a ``splitting'', i.e. with a choice of base change functors $u^* : X_t \longrightarrow X_s$ (for $t \xrightarrow{u} s$ in $S$) \texttt{[rather} $u^* : X_s \longrightarrow X_t$\texttt{]}, with moreover \textit{strict transitivity}
\[
(vu)^* = u^* v^* \qquad\qquad\qquad \text{if}\ \ t \xrightarrow{u} s \xrightarrow{v} r\ \ \text{in}\ S
\]

\noindent\msp{page 4}

\end{document}
